\documentclass[10pt,a4paper]{amsart}
\usepackage{amsmath,amssymb,amsthm,graphicx}
% Retained minimal preamble for the Stacks Project chapter more-morphisms.tex (More on Morphisms).
% Verbatim excerpt of the chapter preamble: Xy-pic setup (native preamble lines 14-18) and the native
% theorem declarations with the mathematical macros (native preamble lines 154-217). Copied mechanically;
% nothing added, removed or redefined. The wrapper class and the amsmath/amssymb/amsthm/graphicx base
% packages are supplied by the assembler, so the class/externaldocument part of the native preamble is omitted.
% For commutative diagrams we use Xy-pic
\usepackage[all]{xy}

% We use 2cell for 2-commutative diagrams.
\xyoption{2cell}
% Theorem environments.
%
\theoremstyle{plain}
\newtheorem{theorem}[subsection]{Theorem}
\newtheorem{proposition}[subsection]{Proposition}
\newtheorem{lemma}[subsection]{Lemma}

\theoremstyle{definition}
\newtheorem{definition}[subsection]{Definition}
\newtheorem{example}[subsection]{Example}
\newtheorem{exercise}[subsection]{Exercise}
\newtheorem{situation}[subsection]{Situation}

\theoremstyle{remark}
\newtheorem{remark}[subsection]{Remark}
\newtheorem{remarks}[subsection]{Remarks}

\numberwithin{equation}{subsection}

% Macros
%
\def\lim{\mathop{\mathrm{lim}}\nolimits}
\def\colim{\mathop{\mathrm{colim}}\nolimits}
\def\Spec{\mathop{\mathrm{Spec}}}
\def\Hom{\mathop{\mathrm{Hom}}\nolimits}
\def\Ext{\mathop{\mathrm{Ext}}\nolimits}
\def\SheafHom{\mathop{\mathcal{H}\!\mathit{om}}\nolimits}
\def\SheafExt{\mathop{\mathcal{E}\!\mathit{xt}}\nolimits}
\def\Sch{\mathit{Sch}}
\def\Mor{\mathop{\mathrm{Mor}}\nolimits}
\def\Ob{\mathop{\mathrm{Ob}}\nolimits}
\def\Sh{\mathop{\mathit{Sh}}\nolimits}
\def\NL{\mathop{N\!L}\nolimits}
\def\CH{\mathop{\mathrm{CH}}\nolimits}
\def\proetale{{pro\text{-}\acute{e}tale}}
\def\etale{{\acute{e}tale}}
\def\QCoh{\mathit{QCoh}}
\def\Ker{\mathop{\mathrm{Ker}}}
\def\Im{\mathop{\mathrm{Im}}}
\def\Coker{\mathop{\mathrm{Coker}}}
\def\Coim{\mathop{\mathrm{Coim}}}

% Boxtimes
%
\DeclareMathSymbol{\boxtimes}{\mathbin}{AMSa}{"02}

%
% Macros for moduli stacks/spaces
%
\def\QCohstack{\mathcal{QC}\!\mathit{oh}}
\def\Cohstack{\mathcal{C}\!\mathit{oh}}
\def\Spacesstack{\mathcal{S}\!\mathit{paces}}
\def\Quotfunctor{\mathrm{Quot}}
\def\Hilbfunctor{\mathrm{Hilb}}
\def\Curvesstack{\mathcal{C}\!\mathit{urves}}
\def\Polarizedstack{\mathcal{P}\!\mathit{olarized}}
\def\Complexesstack{\mathcal{C}\!\mathit{omplexes}}
% \Pic is the operator that assigns to X its picard group, usage \Pic(X)
% \Picardstack_{X/B} denotes the Picard stack of X over B
% \Picardfunctor_{X/B} denotes the Picard functor of X over B
\def\Pic{\mathop{\mathrm{Pic}}\nolimits}
\def\Picardstack{\mathcal{P}\!\mathit{ic}}
\def\Picardfunctor{\mathrm{Pic}}
\def\Deformationcategory{\mathcal{D}\!\mathit{ef}}
\usepackage{cleveref}
\crefname{originallemma}{original lemma}{original lemmas}
\Crefname{originallemma}{Original lemma}{Original lemmas}
\makeatletter
\newlabel{morphisms-lemma-differentials-affine}{{01UT}{source}{}{sourcepointer-1}{}}
\newlabel{morphisms-lemma-differentials-affine@cref}{{[originallemma][0][]01UT}{[1][0][]source}{}{}{}}
\makeatother
\crefname{originaldefinition}{original definition}{original definitions}
\Crefname{originaldefinition}{Original definition}{Original definitions}
\makeatletter
\newlabel{definition-formally-unramified}{{02H8}{source}{}{sourcepointer-2}{}}
\newlabel{definition-formally-unramified@cref}{{[originaldefinition][0][]02H8}{[1][0][]source}{}{}{}}
\makeatother
\title{A morphism is formally unramified exactly when its sheaf of differentials vanishes}
\author{The Stacks Project Authors}
\date{stacks-project:more-morphisms.tex:local-snapshot-20261011; GNU FDL 1.2 or later}
\begin{document}
\maketitle
\noindent\textit{Editorial note.} All delivered mathematical text is the author's own, copied from the pinned Stacks Project chapter file more-morphisms.tex; no statement or proof step was rewritten, reordered or removed. The author's slip 'defined by a an ideal' is reproduced verbatim.\par\medskip
\noindent\textit{Editorial note.} The preamble is the chapter's own retained minimal preamble (sources/stacks-more-morphisms/retained-preamble.tex), a verbatim excerpt of the native Xy-pic setup and of the native theorem declarations and mathematical macros; no macro was invented or redefined.\par\medskip
\noindent\textit{Editorial note.} The same-chapter reference is located at its \textbackslash{}label line inside the pinned chapter file; the cross-chapter pointer is located in the project's own tag list staged in the snapshot (sources/stacks-chow/tags.txt), at the verbatim line carrying its canonical tag.\par\medskip
\noindent\textit{Allowed standard prerequisites (editorial).}\begin{itemize}
\item Formally unramified morphism, tested against first order thickenings. Same chapter, Definition environment labelled definition-formally-unramified (source line 798, Stacks tag 02H8); invoked at proof line 934. Supplies the two-morphism lifting condition that the converse direction has to verify; the verification itself is retained in the counted span.
\item The conormal sequence of the diagonal and the identification of the sheaf of differentials with J/J\textasciicircum{}2. Morphisms, Lemma labelled morphisms-lemma-differentials-affine (canonical Stacks tag 01UT); the author says at proof lines 903-905 that the diagrammatic set-up is recalled from its proof. Supplies the open W containing the image of the diagonal, the ideal J of the diagonal, the square-zero ideal J\textasciicircum{}2 and the identification Omega\_\{X/S\} = J/J\textasciicircum{}2 used throughout.
\end{itemize}\medskip
\ifdefined\hypertarget\hypertarget{sourcepointer-1}{}\fi
\noindent\textit{Source reference.} Original lemma 01UT: sources/stacks-chow/tags.txt, source lines 2422--2422. Cross-chapter standard prerequisite invoked at proof line 905 as the place where the diagonal set-up is recalled from; the author reproduces the needed facts (W, J, J\textasciicircum{}2, Omega\_\{X/S\} = J/J\textasciicircum{}2) inside the counted span.\par\medskip
\ifdefined\hypertarget\hypertarget{sourcepointer-2}{}\fi
\noindent\textit{Source reference.} Original definition 02H8: sources/stacks-more-morphisms/pinned-original-source.tex, source lines 798--798. Same-chapter pointer invoked at proof line 934 for the lifting condition that the converse direction verifies; that verification is retained verbatim in the counted span.\par\medskip
\begin{lemma}
\label{lemma-formally-unramified-differentials}
Let $f : X \to S$ be a morphism of schemes.
Then $f$ is formally unramified if and only if $\Omega_{X/S} = 0$.
\end{lemma}

\begin{proof}
We recall some of the arguments of the proof of
Morphisms, Lemma \ref{morphisms-lemma-differentials-affine}.
Let $W \subset X \times_S X$ be an open such that
$\Delta : X \to X \times_S X$ induces a closed immersion into $W$.
Let $\mathcal{J} \subset \mathcal{O}_W$ be the ideal sheaf of this
closed immersion. Let $X' \subset W$ be the closed subscheme
defined by the quasi-coherent sheaf of ideals $\mathcal{J}^2$.
Consider the two morphisms $p_1, p_2 : X' \to X$ induced by
the two projections $X \times_S X \to X$.
Note that $p_1$ and $p_2$ agree when composed with $\Delta : X \to X'$
and that $X \to X'$ is a closed immersion defined by a an ideal
whose square is zero. Moreover there is a short exact sequence
$$
0 \to \mathcal{J}/\mathcal{J}^2 \to \mathcal{O}_{X'} \to \mathcal{O}_X \to 0
$$
and $\Omega_{X/S} = \mathcal{J}/\mathcal{J}^2$. Moreover,
$\mathcal{J}/\mathcal{J}^2$ is generated by the local
sections $p_1^\sharp(f) - p_2^\sharp(f)$ for $f$ a local section of
$\mathcal{O}_X$.

\medskip\noindent
Suppose that $f : X \to S$ is formally unramified.
By assumption this means that $p_1 = p_2$ when restricted to any
affine open $T' \subset X'$. Hence $p_1 = p_2$. By what was said above
we conclude that $\Omega_{X/S} = \mathcal{J}/\mathcal{J}^2 = 0$.

\medskip\noindent
Conversely, suppose that $\Omega_{X/S} = 0$. Then $X' = X$. Take any pair
of morphisms $f'_1, f'_2 : T' \to X$ fitting as dotted arrows in
the diagram of
Definition \ref{definition-formally-unramified}.
This gives a morphism $(f'_1, f'_2) : T' \to X \times_S X$.
Since $f'_1|_T = f'_2|_T$ and $|T| =|T'|$ we see that the image of $T'$
under $(f'_1, f'_2)$ is contained in the open $W$ chosen above. Since
$(f'_1, f'_2)(T) \subset \Delta(X)$ and since $T$ is defined by an ideal
of square zero in $T'$ we see that $(f'_1, f'_2)$ factors through $X'$.
As $X' = X$ we conclude $f_1' = f'_2$ as desired.
\end{proof}

\end{document}
